\documentclass[10pt,a4paper]{amsart}
\usepackage[margin=19mm]{geometry}
\usepackage{amsmath,amssymb,mathtools}
\usepackage[scr=rsfs,cal=euler]{mathalfa}
\usepackage{xfrac}
\usepackage[unicode,hidelinks,bookmarks=false]{hyperref}
\newcommand{\ie}{i.e.\ }
\newcommand{\cf}{cf.\ }
\newcommand{\resp}{respectively\ }
\newcommand{\Subsection}{\subsection}
\providecommand{\nobreakdash}{\nobreak\hbox{-}}
\setlength{\emergencystretch}{2em}
\multlinegap0pt
\usepackage{xcolor}
\usepackage{amssymb}
\newtheorem{thm}{Theorem}
\title{On the Capacity of Zero-Drift First Arrival Position Channels in Diffusive Molecular Communication \iffalse {\footnotesize \textsuperscript{*}Note: Sub-titles are not captured in Xplore and should not be used} \thanks{Identify applicable funding agency here. If none, delete this.} \fi}
\author{Yen-Chi Lee, Min-Hsiu Hsieh}
\date{Source: arXiv:2201.11383v3 (CC BY 4.0)}
\begin{document}
\maketitle

\noindent\emph{Source and licence.} On the Capacity of Zero-Drift First Arrival Position Channels in Diffusive Molecular Communication \iffalse {\footnotesize \textsuperscript{*}Note: Sub-titles are not captured in Xplore and should not be used} \thanks{Identify applicable funding agency here. If none, delete this.} \fi. Version arXiv:2201.11383v3 (CC BY 4.0), distributed by arXiv under the Creative Commons Attribution 4.0 International licence, http://creativecommons.org/licenses/by/4.0/, as recorded in the arXiv licence metadata for this exact version and shown on the abstract page https://arxiv.org/abs/2201.11383v3. Full licence notice: \texttt{licenses/arxiv-2201.11383-CC-BY-4.0.txt}.\par\smallskip

\noindent\textit{Editorial scope.} Counted unit: the article's thm thm:T2 (source0.tex lines 737--752). The article's complete original proof of it is delivered with it (source0.tex lines 754--791, one \texttt{proof} environment of 38 lines). No mathematics is written, repaired or re-derived: the statement and the proof are the article's own lines, in the article's own notation, and no retained line is substituted or re-flowed.  Retained uncounted as the source-local context the proof needs: lines 702--705.  Nothing was omitted from any retained span, so the package carries no omission record.  No retained line contains a citation, so the wrapper carries no bibliography and no bibliography entry.  No retained line contains a figure environment, an image-inclusion command or drawing code, so no image is delivered and the package declares no asset dependency.  Editorial formatting only, in addition: an AMS \texttt{amsart} preamble that restates the article's own declarations and macros used by the retained lines, namely the environments \texttt{thm} on separate counters, together with the standard packages those lines need.  The retained environments print in this document as Theorem 1 in order of appearance, whereas the article prints the same items with the numbering of its own full document.  The complete native source of the article as distributed in the arXiv e-print for this exact version is bundled unchanged as \texttt{sources/121276/source0.tex}.
\begin{equation}
    \mathbf{Y}=\mathbf{X}+\mathbf{N},
\label{eq:45}
\end{equation}

\begin{thm}
The capacity of 3D FAP channel \eqref{eq:45}  is
\begin{equation}
    C_{\rm{3D,FAP}}=C(A,\lambda)=2\ln\left(\dfrac{A}{\lambda}\right)
\end{equation}
under the output logarithmic constraint
$
\mathbf{Y}\in \mathcal{D}_{{\rm bi}}(A)
$
for some 
prescribed dispersion level $A$, where $A \geq \lambda$. In addition, the corresponding capacity achieving output distribution is
\begin{equation}
\mathbf{Y}^* \sim {\rm Cauchy}_2(0,{\rm diag}(A^2,A^2)).
\end{equation}
\label{thm:T2}
\end{thm}

\begin{proof}
%Next we prove Theorem~\ref{thm:T2}. 
For the additive channel we are considering, the mutual information can be written as
\begin{equation}
I(\mathbf{X};\mathbf{Y})=h(\mathbf{Y})-h(\mathbf{N}).
\end{equation}
The problem of finding the channel capacity turns out to be an optimization problem:
$
C(A,\lambda)=\sup_{\mathbf{Y}\in\mathcal{D}_{\text{bi}}(A)} I(\mathbf{X};\mathbf{Y}).
$
We have
\begin{align}
    C(A,\lambda)
    &=\sup_{\mathbf{Y} \in \mathcal{D}_{\text{bi}}(A)} \Big\{h(\mathbf{Y} ) - h(\mathbf{N} )\Big\}\\
    &=\Big\{\sup_{\mathbf{Y} \in \mathcal{D}_{\text{bi}}(A)} h(\mathbf{Y} )\Big\} - \ln(2\pi e^3\lambda^2), \label{eq:54}
\end{align}
where
\begin{align}
    \sup_{\mathbf{Y} \in \mathcal{D}_{\text{bi}}(A)} h(\mathbf{Y})
    &=
    \sup_{k\in[\lambda,A]}
    \Big\{ \sup_{\mathbf{Y}\in \mathcal{D}_{k,\text{bi}}} h(\mathbf{Y}) \Big\}\label{eq:55}\\
    &=
    \sup_{k\in[\lambda,A]}
    \ln(2\pi e^3 k^2) \label{eq:56}\\
    &= \ln(2\pi e^3 A^2). \label{eq:57}
\end{align}
Combining equations \eqref{eq:54} and \eqref{eq:57} yields
\begin{align}
\begin{split}
C(A,\lambda)
=&\ \ln(2\pi e^3 A^2) - \ln(2\pi e^3 \lambda^2)\\
=&\ \ln\left(\dfrac{A^2}{\lambda^2}\right)
=2\ln\left(\dfrac{A}{\lambda}\right).
\end{split}
\end{align}
Notice that the equalities \eqref{eq:55}-\eqref{eq:57} occur if and only if $\mathbf{Y}$ distributes as ${\rm Cauchy}_2(0,\text{diag}(A^2,A^2))$. Hence, we have proved Theorem~\ref{thm:T2}.
\end{proof}

\end{document}
